\begin{table}[t]
\centering
\caption{Numerical data underlying Fig.~5. Operational metrics for Full ARMG relative to the stateless self-correction baseline.}
\label{tab:fig5_validation}
\footnotesize
\begin{tabular}{lccccc}
\toprule
\textbf{Evaluation Metric} & \textbf{Mode 2} & \textbf{Mode 4} & \textbf{Figure Delta} & \textbf{Delta Type} & \textbf{Unit} \\
\midrule
Mean Repair Iterations & 0.28 & 0.37 & +33.33\% & Relative & retries/query \\
Mean Token Expenditure & 503.72 & 602.85 & +19.68\% & Relative & tokens/query \\
PostgreSQL Execution Success & 96.00\% & 96.00\% & +0.00~pp & Percentage points & pp \\
End-to-End Latency & 6,441.56 & 9,000.59 & +39.73\% & Relative & ms \\
Relational Execution Accuracy & 68.00\% & 68.00\% & +0.00~pp & Percentage points & pp \\
\bottomrule
\multicolumn{6}{p{8.3cm}}{\scriptsize \textsuperscript{*}Primary source: \texttt{benchmark/seed*/benchmark\_results.csv}, Table~II, Table~VIII, and \texttt{manuscript/figures/source/fig5\_tradeoff.py}. Mean values represent 75 queries evaluated per mode across seeds 42, 123, and 999. In strict accordance with IEEE scientific reporting conventions, efficiency shifts for rate metrics (retries, tokens, latency) are reported as relative percentage changes (\%), whereas probability metrics bounded in $[0, 100]$ (execution success, semantic accuracy) are reported as arithmetic percentage-point differences (pp).} \\
\end{tabular}
\end{table}
